%% !TEX root = ../main.tex
%% ================================================================
%%  Chapter 1 - Introduction
%%  Skeleton + guidance. Replace the italic notes with your prose.
%%  Worked LaTeX examples live in Appendix B.
%% ================================================================
\chapter{Introduction}
\label{ch:introduction}

\section{Background and Motivation}
\label{sec:motivation}

\emph{(Set the scene. What is the real-world problem, and why does it matter?
Give the reader the context they need in a few paragraphs, and cite the key
prior work~\cite{smith2024example}. End the section by narrowing from the broad
context down to the specific gap your thesis addresses.)}

\section{Problem Statement}
\label{sec:problem}

\emph{(State precisely the problem you tackle. What is unknown, missing, or
unsatisfactory in current approaches? Make the gap concrete and specific.)}

\section{Research Questions and Objectives}
\label{sec:objectives}

\emph{(State the objective of the thesis in one sentence, then list the concrete
questions it answers.)}

\begin{enumerate}[itemsep=0.3em]
  \item \emph{(First research question.)}
  \item \emph{(Second research question.)}
  \item \emph{(Third research question.)}
\end{enumerate}

\section{Contributions}
\label{sec:contributions}

\emph{(List what is new in your work. Each bullet should be a concrete,
defensible claim that the rest of the thesis supports.)}

\begin{itemize}[itemsep=0.3em]
  \item \emph{(First contribution.)}
  \item \emph{(Second contribution.)}
  \item \emph{(Third contribution.)}
\end{itemize}

\section{Thesis Outline}
\label{sec:outline}

\emph{(One short paragraph mapping the rest of the document.)}
\Cref{ch:background} reviews the related work; \Cref{ch:data} describes the data
and materials; \Cref{ch:methodology} presents the method; \Cref{ch:results}
reports and discusses the results; and \Cref{ch:conclusion} concludes and
outlines future work.
